\chapter{Introducción a CVE, CVSS y las principales bases de datos}

El programa \textbf{CVE} (\textit{Common Vulnerabilities and Exposures}) es un catálogo público y estandarizado de vulnerabilidades de seguridad informática, gestionado por la corporación \textbf{MITRE} desde 1999 con patrocinio de la agencia estadounidense \textit{Cybersecurity and Infrastructure Security Agency} (CISA). Su objetivo es asignar a cada vulnerabilidad un identificador único (de la forma \texttt{CVE-AÑO-NÚMERO}) que permita a fabricantes, investigadores y administradores referirse a un mismo fallo sin ambigüedad. La asignación de estos identificadores no es centralizada: MITRE delega esa potestad en las \textbf{CNA} (\textit{CVE Numbering Authorities}), que son organizaciones autorizadas --- típicamente fabricantes de software (Microsoft, Apple, Google, Red Hat, etc.), CERTs nacionales o programas de \textit{bug bounty} como \textit{Zero Day Initiative} --- encargadas de reservar, documentar y publicar las CVE dentro de su propio ámbito.

El \textbf{CVSS} (\textit{Common Vulnerability Scoring System}), mantenido por el foro internacional \textbf{FIRST.org}, es la métrica de facto para cuantificar la gravedad de cada vulnerabilidad. La puntuación final (entre 0.0 y 10.0) se deriva de un \textit{vector} con varias métricas base: vector de ataque (AV), complejidad (AC), privilegios necesarios (PR), interacción del usuario (UI), alcance (S) y los impactos sobre confidencialidad, integridad y disponibilidad (C, I, A). El estándar establece cinco niveles cualitativos (\textit{None, Low, Medium, High, Critical}), considerando \textit{Critical} las puntuaciones $\geq 9.0$. La versión vigente para la mayoría de registros actuales es \textbf{CVSS v3.1}, si bien desde 2023 coexiste con \textbf{CVSS v4.0}.

Existen tres bases de datos complementarias que conviene distinguir. \textbf{CVE.org} es el portal oficial de MITRE y contiene los registros tal como los publican las CNA, es decir, la fuente primaria de cada vulnerabilidad. La \textbf{NVD} (\textit{National Vulnerability Database}, \texttt{nvd.nist.gov}) la mantiene el \textbf{NIST} estadounidense y \textit{enriquece} esos registros añadiendo el cálculo CVSS oficial, la traducción a CPE (\textit{Common Platform Enumeration}) y referencias a CWE; es la base de datos de consulta más utilizada en la práctica profesional. Por último, \textbf{FIRST.org} no es un repositorio de CVE, sino la organización custodia del estándar CVSS: publica la especificación, el calculador oficial y la metodología de puntuación. En esta práctica se ha trabajado fundamentalmente con la NVD, complementada con consultas puntuales a CVE.org y a la propia documentación de los fabricantes afectados.
